%% Thesis language: english or polish
%% Thesis type: master or engineer
\documentclass[polish,engineer]{polsl-msth}
%% set document encoding
\usepackage[utf8]{inputenc}

%% thesis author
\author{Jan Kowalski}

%% thesis supervisor
\supervisor{Jan Kowalski}

%% thesis consultant, optional
%% \consultant{Jan Kowalski}

%% Thesis title
\title{Przykład wykorzystania klasy polsl-msth}
\begin{document}
\frontmatter
\maketitle
\tableofcontents
\listoftables
\listoffigures
\mainmatter
\chapter{Wprowadzenie}
Dokument ten opisuje wykorzystanie klasy polsl-msth do przygotowania
pracy magisterskiej na Wydziale Automatyki, Elektroniki i~Informatyki
Politechniki Śląskiej. Sam dokument jest
też zgodny z~szablonem pracy magisterskiej.

W~rozdziale \ref{zalecenia} znajdują się zalecenia odnośnie formatowania
pracy.

\chapter{Zalecenia\label{zalecenia}}
Pierwszy akapit rozdziału, sekcji lub podsekcji powinien rozpoczynać
się bez wcięcia. Wszystkie kolejne powinny mieć wcięcie skonfigurowane
w~klasie.

Opis rysunków powinien znajdować się pod rysunkiem jak w~przypadku
rysunku \ref{pi}.
\begin{figure}[hbtp]
\centerline{\includegraphics[width=1in]{pi}}
\caption{Logo Politechniki Śląskiej.\label{pi}}
\end{figure}

Opis tabel powinien znajdować się nad tabelą jak w~przypadku
tabeli \ref{tab}.
\begin{table}[hbtp]
\caption{Przykładowa tabela.\label{tab}}
\centerline{\begin{tabular}{|c|c|}
\hline
1 & 2\\
\hline
3 & 4\\
\hline
\end{tabular}}
\end{table}

Równania powinny być numerowane jak równanie \ref{entropy}.
\begin{equation}
H(X) = -\sum_{i=0}^n P(x_i) \log_2 P(x_i).\label{entropy}
\end{equation}
\end{document}
